\section{Conclusion}\label{sec:conclusion}

We studied a two-stage cascade for the detection of fraudulent short text messages: a calibrated linear model over character n-grams that answers when its confidence reaches a threshold, and a DistilBERT classifier for the remaining messages. The threshold is the solution of an explicit selection problem on validation data, and obfuscation is modelled as a random string transformation with one family of operations for training and another held out for testing. The evaluation used deduplicated corpora, a protocol frozen before the test data were read, \val{n-seeds} seeds, and paired tests with Holm correction.

The results support three conclusions. First, a calibrated character-level linear model is a strong baseline for this task. On the clean three-class test set it reached the highest macro F1 (\val{svmchar-clean-f1}), and the cascade (\val{cascade-clean-f1}) did not differ from it significantly; across all conditions the cascade was significantly better than the best baseline only twice, each time by one of the two tests. Second, robustness to character obfuscation comes from the character level. Under the strongest held-out obfuscation, the re-implemented published methods fell to between \val{pub-ho20-min} and \val{pub-ho20-max} in their own settings, while the cascade and the character model alone stayed at or above \val{cascade-shim-ho20}; augmentation made the transformer more robust but did not improve the cascade. Third, the cascade is fast when its first stage is confident: in the three-class task it sent \val{lat-cascade-routed-pct}\% of the test messages to the transformer, with a median latency of \val{lat-cascade-p50}~ms on one CPU thread, and it had the lowest error cost of all models under strong held-out obfuscation. When the first stage is often not confident, as with little training data or after a change of language, a larger share of the messages reaches the transformer (up to \val{tr-ft-share}\% in the transfer experiment), and the cascade inherits its cost and its weaknesses.

The duplicate analysis showed that the handling of duplicates and conflicting labels changes cross-validated macro F1 by up to \val{infl-max-drop} points, in a direction that depends on the task, which is more than most differences between models. Results on public corpora of short messages should therefore be reported with an explicit deduplication procedure and with significance tests.

The case study for Kazakhstan showed that models trained on English messages miss most fraudulent messages written in Russian or Kazakh, that machine translation helps only partly and is too slow for real-time use, and that a small synthetic adaptation set raises recall at the price of false alarms on legitimate notifications. Future work should collect real labelled messages from such settings, evaluate obfuscation chosen by an adversary rather than at random, use thresholds that account for the cost of errors, separate the effects of duplicates and conflicting labels, and test a multilingual first stage.
